\section{Inference}
\label{sec:inference}

Earlier drafts of this paper reported every improvement against one shared
comparator, applied no multiplicity control, and left the central model-class
difference untested. This section repairs all three. It also reports what the
repairs cost: two arms lose significance, one claimed effect does not survive
family-wise correction, and one that the paper leaned on turns out never to
have been resolved at all.

\subsection{Multiplicity Across the Battery}
\label{sec:multiplicity_results}

Ninety-seven arms are compared against the shared incumbent, and the paper
quotes the minimum in two places. Table-wide correction is therefore
mandatory, not optional.

Of the 97 tests, 93 are significant at an unadjusted $p < 0.05$. Applying
Holm--Bonferroni across the full family leaves \textbf{91}; Benjamini--Hochberg
at the same level leaves 93. The two arms that lose significance under Holm are
\texttt{reclasso/sentiment} ($p = 0.036 \to 0.215$) and the base-5 ladder arm
\texttt{ols\_har b5 cap960} ($p = 0.037 \to 0.215$).

The conclusion is favourable to the paper and worth stating plainly: the
information and model-class effects are not artifacts of multiple testing. A
$t$ of $-19$ against a family of 97 survives any correction one might apply,
and even the weakest bucket effect we claim---\texttt{vol\_demand} at
$t = -4.1$---clears Holm comfortably. What does not survive is the marginal
sentiment arm, which we no longer claim.

\subsection{The Pairwise Model-Class Test}
\label{sec:pairwise}

The paper's model-class conclusion compares two arms that were each tested
only against a third. That is not the test the claim requires, and the
per-bar loss series needed for a direct Diebold--Mariano test between them are
not retained. A bound is available without them.

For two arms $A$ and $B$ scored against a common incumbent, the incumbent
cancels from the pairwise differential:
\begin{equation}
  d_A - d_B \;=\; (L_A - L_{\text{inc}}) - (L_B - L_{\text{inc}}) \;=\; L_A - L_B ,
  \label{eq:pairwise_cancel}
\end{equation}
so the quantity we want is exactly the difference of two reported means. Its
variance is
\begin{equation}
  \operatorname{Var}\!\left(\overline{L_A - L_B}\right)
  \;=\; se_A^2 + se_B^2 - 2\rho\, se_A se_B ,
  \label{eq:pairwise_var}
\end{equation}
with $se_A$, $se_B$ recoverable from each arm's reported mean differential and
HAC $t$ statistic. Taking the worst case $\rho = -1$ gives
$se \le se_A + se_B$, hence a \emph{lower bound} on $|t|$ and an \emph{upper
bound} on $p$ that assume nothing whatever about the dependence between the two
arms. If the bound clears the critical value the difference is established; if
it does not, the comparison is inconclusive rather than negative.

\begin{table}[H]
\centering
\caption{Pairwise model-class inference on the frozen battery, from the
reported HAC standard errors alone. For arms sharing an incumbent the
pairwise differential is exactly $L_A-L_B$, whose standard error is bounded
above by $se_A+se_B$ regardless of the correlation between the two arms'
differentials; the resulting $|t|$ is therefore a \emph{lower bound} and the
$p$ an \emph{upper bound}. A tighter bound assuming the two arms'
improvements co-move ($\rho\ge 0$) is reported in the text; it does not
change any verdict. Holm adjustment is applied across the nine matched
comparisons. Negative $t$ favours the tree.}
\label{tab:pairwise_bounds}
\begin{tabular}{lrrrrl}
\toprule
Feature set & $\Delta$QLIKE & $|t|\ge$ (free) & $p\le$ (free) & $p_{\text{Holm}}\le$ & Verdict \\
\midrule
\texttt{all\_features} & -0.00184 & 2.19 & 0.029 & 0.029 & tree better \\
\texttt{moments} & -0.00269 & 3.82 & $1.4\times 10^{-4}$ & $4.1\times 10^{-4}$ & tree better \\
\texttt{liquidity} & -0.00281 & 4.59 & $4.5\times 10^{-6}$ & $2.3\times 10^{-5}$ & tree better \\
\texttt{market\_vw} & -0.00327 & 5.91 & $3.4\times 10^{-9}$ & $3.1\times 10^{-8}$ & tree better \\
\texttt{vol\_demand} & -0.00358 & 5.71 & $1.1\times 10^{-8}$ & $9.1\times 10^{-8}$ & tree better \\
\texttt{implied\_vol} & -0.00252 & 3.73 & $1.9\times 10^{-4}$ & $4.1\times 10^{-4}$ & tree better \\
\texttt{sentiment} & -0.00310 & 5.27 & $1.4\times 10^{-7}$ & $8.3\times 10^{-7}$ & tree better \\
\texttt{market\_ew} & -0.00281 & 4.17 & $3.0\times 10^{-5}$ & $1.2\times 10^{-4}$ & tree better \\
\texttt{baseline} & -0.00327 & 5.54 & $3.0\times 10^{-8}$ & $2.1\times 10^{-7}$ & tree better \\
\bottomrule
\end{tabular}
\end{table}


\textbf{The model-class ordering is established.} All nine matched comparisons
resolve in the tree's favour after Holm adjustment, under a bound that makes no
assumption about the dependence between arms. The headline pair---best tree
$0.12604$ against best linear $0.12788$---is the weakest of the nine, at
$|t| \ge 2.19$ and $p \le 0.029$. Under the tighter bound that holds when the
two arms' improvements co-move ($\rho \ge 0$), the same pair gives
$|t| \ge 3.00$, $p \le 0.0027$; we quote the assumption-free figure.

Two qualifications. First, these are bounds: the true statistics are larger in
magnitude, so the evidence is stronger than the table shows, but by an
unmeasured amount. Second, the headline pair is only just resolved. A referee
would be right to note that the paper's most-quoted number, $0.00184$, is the
least secure of the nine, and that the case for the model-class ordering rests
more on its consistency across feature sets than on any single comparison.

\paragraph{A distribution-free check.} The tree beats the linear arm on all
nine feature sets. Under a null of no model-class effect, the exact binomial
sign test gives $p = 0.0039$. The nine feature sets are nested and share a
panel, so this is anti-conservative and we report it as a consistency check
rather than as an independent test.

\subsection{Per-Bar Inference on the Campaign Tile}
\label{sec:tile_inference}

Several claims in this paper are inherited from campaign experiments whose
per-bar predictions \emph{are} retained, on a common tile of $n = 2{,}189$ bars
with identical targets across 161 arms. That tile permits genuine
Diebold--Mariano tests and genuine model confidence sets rather than bounds.

\paragraph{Panel discipline.} The tile is not the frozen battery. Its levels are
not comparable to any level elsewhere in this paper, and its sample is two
orders of magnitude smaller, so it has far less power. What it can settle is
whether each \emph{claim} survives per-bar inference at all. We report it
separately for that reason and never mix its numbers with battery numbers.

\begin{table}[H]
\centering
\caption{Per-bar Diebold--Mariano tests on the campaign tile ($n = 2,189$ bars, 161 arms, identical targets). HAC (Newey--West) standard errors with the
Harvey--Leybourne--Newbold small-sample correction; Holm adjustment across
the whole contrast family. Negative $t$ favours arm A. \textbf{Panel note:}
this is a campaign tile, not the 218{,}934-bar frozen battery; no level here
is comparable to a battery level, and what the tile tests is whether each
claim survives per-bar inference at all.}
\label{tab:tile_dm}
\footnotesize
\begin{tabular}{lrrrl}
\toprule
Contrast (A vs B) & $\Delta$QLIKE & $t$ & $p_{\text{Holm}}$ & Verdict \\
\midrule
trees direct: xgb vs incumbent & +0.10163 & 6.96 & $8.7\times 10^{-11}$ & B better \\
local k=8000 vs k=200 & -0.02041 & -6.59 & $9.8\times 10^{-10}$ & A better \\
trees direct: lgbm vs incumbent & +0.07684 & 6.46 & $2.3\times 10^{-9}$ & B better \\
cap256 vs cap32 & -0.00735 & -4.16 & $5.2\times 10^{-4}$ & A better \\
cap256 vs uncapped & -0.00774 & -3.83 & 0.002 & A better \\
local k=16000 vs k=8000 & +0.00294 & 3.74 & 0.003 & B better \\
cap256 vs cap128 & -0.00264 & -3.50 & 0.006 & A better \\
enet mid vs ridge & +0.00523 & 3.15 & 0.020 & B better \\
rank vs divide & +0.00718 & 2.73 & 0.070 & unresolved \\
enet heavy vs ridge & +0.00573 & 2.20 & 0.277 & unresolved \\
Fourier clock vs session gates & -0.00367 & -1.64 & 0.908 & unresolved \\
PCovR-50 vs PCovR-25 & -0.00213 & -1.56 & 0.954 & unresolved \\
trees residual: xgb vs ridge cadence & +0.00321 & 1.11 & 1.000 & unresolved \\
trees residual: lgbm vs ridge cadence & +0.00308 & 1.06 & 1.000 & unresolved \\
PLS reg vs PCovR & +0.00054 & 0.68 & 1.000 & unresolved \\
enet light vs ridge & -0.00012 & -0.65 & 1.000 & unresolved \\
tie: rawc stack vs champion & +0.00063 & 0.30 & 1.000 & unresolved \\
tie: prune32+corr vs champion & +0.00014 & 0.06 & 1.000 & unresolved \\
tie: legchamp+clock vs champion & +0.00001 & 0.00 & 1.000 & unresolved \\
\bottomrule
\end{tabular}
\end{table}


Four claims are confirmed at Holm-adjusted significance across the
nineteen-contrast family.

\begin{itemize}
  \item \textbf{Trees fail as direct forecasters on the kernel dictionary}, and
    fail badly: XGBoost $+0.102$ ($t = 6.96$) and LightGBM $+0.077$
    ($t = 6.46$) against the incumbent. The extrapolation clamp reported in
    Section~\ref{sec:saturation} is real and is the largest single effect in
    the paper.
  \item \textbf{Heavy local averaging pays}, decisively: $k = 8000$ beats
    $k = 200$ by $0.0204$ ($t = -6.59$).
  \item \textbf{Kernel truncation is a genuine interior optimum.} Cap-256 beats
    cap-32 ($-0.0074$), beats no truncation at all ($-0.0077$), and beats
    cap-128 ($-0.0026$), all Holm-significant. Truncating the kernel is not a
    computational convenience; removing it costs real accuracy.
  \item \textbf{Selection that bites hurts.} Elastic net at $\alpha = 3\times
    10^{-4}$ is $+0.0052$ worse than ridge ($t = 3.15$), while light $\ell_1$
    is indistinguishable from ridge. This is the sharpest available form of
    ``shrink, never select''.
\end{itemize}

Two results refine claims made elsewhere in the paper.

\begin{itemize}
  \item \textbf{Averaging has an interior optimum too.} Moving from $k = 8000$
    to $k = 16000$ is $+0.0029$ \emph{worse} ($t = 3.74$, Holm-significant). The
    correct statement is not ``more averaging is better'' but that the optimum
    sits at a large but finite neighbourhood; Section~\ref{sec:spectral_knn}
    is amended accordingly.
  \item \textbf{The parsimony ties are ties.} The three compressed
    representations are statistically indistinguishable from the wide champion:
    $p = 0.998$, $0.951$ and $0.763$ respectively. A tie is not evidence of
    equality, but with these $p$-values the burden plainly lies with anyone
    claiming the wide model does something the compressed ones do not.
\end{itemize}

\subsubsection{Three claims that do not survive}
\label{sec:not_surviving}

We report these at the same altitude as the confirmations.

\begin{enumerate}
  \item \textbf{Trees as residual readers are not shown to fail.} LightGBM
    $+0.0031$ ($t = 1.06$, $p = 0.29$) and XGBoost $+0.0032$ ($t = 1.11$,
    $p = 0.27$) against the smooth linear reader. The sign is as claimed and
    the point estimates match earlier reports, but neither is distinguishable
    from zero. The claim that ``trees fail in both seats'' holds for the direct
    seat only.
  \item \textbf{Fourier clock modulation does not beat session gates.}
    $-0.0037$ with $t = -1.64$, $p = 0.10$ unadjusted and $0.91$ after Holm.
    An internal note recorded gates as losing on both tiles; on this tile that
    is a point estimate, not a result.
  \item \textbf{Rank versus divide is not resolved.} $+0.0072$ with
    $t = 2.73$, $p = 0.0064$ unadjusted but $0.071$ after Holm across the
    contrast family. An earlier note recorded $p = 0.001$ for this comparison;
    under HAC standard errors with a small-sample correction and family-wise
    adjustment it falls short at the 5\% level.
\end{enumerate}

\subsection{Model Confidence Sets}
\label{sec:mcs_results}

\begin{table}[H]
\centering
\caption{Hansen--Lunde--Nason model confidence sets on the campaign tile ($\alpha = 0.1$, $B = 5,000$ stationary-bootstrap replications, $T_{\max}$ statistic). ``Kept'' counts arms that cannot be excluded from
the set of best models. Survivor counts for the top-twenty family are
unchanged when the bootstrap block length is varied from half to four times
its automatic value.}
\label{tab:mcs}
\begin{tabular}{llrrl}
\toprule
Family & Question & Arms & Kept & Eliminated \\
\midrule
\texttt{penalty} & Penalty: shrinkage vs selection & 4 & 2 & \texttt{pbe\_enet\_0.0003}, \texttt{pbe\_enet\_0.003} \\
\texttt{trees\_vs\_linear} & Model class on the kernel dictionary & 6 & 4 & \texttt{pbt\_xgb\_cad}, \texttt{pbt\_lgbm\_cad} \\
\texttt{local\_k} & Analog neighbourhood size & 6 & 3 & \texttt{i\_loc\_k200}, \texttt{incumbent}, \texttt{i\_loc\_k16000} \\
\texttt{cap\_ladder} & Kernel truncation depth & 6 & 6 & --- \\
\texttt{dim\_reduction} & Dimension reduction & 8 & 8 & --- \\
\texttt{top20} & The twenty best arms & 20 & 20 & --- \\
\texttt{all} & Every tile arm & 161 & 123 & (see text) \\
\bottomrule
\end{tabular}
\end{table}


The MCS is a joint test and correspondingly less powerful than the targeted
contrasts above, which shows in the results.

\paragraph{The discriminating comparisons discriminate.} On the penalty family
the MCS retains ridge and light elastic net and eliminates both heavier
elastic nets ($p = 0.009$, $0.023$). On the model-class family it retains
ridge, both residual-reading trees and the incumbent, and eliminates both
direct-forecasting trees ($p < 0.001$). On the neighbourhood family it retains
$k \in \{1000, 4000, 8000\}$ and eliminates $k = 200$, $k = 16000$ and the
incumbent.

\paragraph{The top of the tile is a wide tie.} All twenty of the best arms
survive, and the result is unchanged when the bootstrap block length is varied
from half to four times its automatic value. Over the full set of 161 arms the
MCS retains 123. With $n = 2{,}189$ this is a statement about power, not about
the models: the tile cannot discriminate among the great majority of arms
tried on it, and any ranking within that set---including the ordering of the
top three---is not supported by the data.

\paragraph{Consequence for the campaign.} The practice of reading a tile
leaderboard as a ranking is not defensible at this sample size. Differences at
the top of the tile are real only if they replicate at scale, which is the
argument for the frozen-battery discipline adopted for the rest of this paper.

\subsection{Auxiliary Tests}
\label{sec:aux_tests}

\paragraph{Ladder monotonicity.} The quadrature prediction of
Section~\ref{sec:powerlaw} is that QLIKE increases monotonically with the
ladder base. Across the six bases the rank correlation is exactly $+1.000$
($p < 10^{-4}$): the ordering is perfect, with no inversions.

\paragraph{Metric disagreement.} The rank correlation between QLIKE and
raw-variance OOS-$R^2$ across the 45 causally tuned arms is $+0.152$ with
$p = 0.32$. Agreement between the two metrics would produce a \emph{negative}
correlation; what we observe is indistinguishable from no relationship at all.
The two metrics do not rank this battery in a related way.

\paragraph{Calibration.} The rank correlation between QLIKE and the
Mincer--Zarnowitz slope is $+0.741$ with $p = 5.8 \times 10^{-9}$. Arms that
score better on QLIKE are more attenuated, and the association is
overwhelmingly significant. Section~\ref{sec:calibration} discusses what this
implies for reading a QLIKE ranking.

\subsection{What Remains Unresolved}
\label{sec:inference_remaining}

\begin{itemize}
  \item A full model confidence set over the 45 causally tuned battery arms
    still requires their per-bar losses, which are not retained. The pairwise
    bounds of Section~\ref{sec:pairwise} settle the model-class question but do
    not produce a confidence set over the battery.
    \pending{re-run the battery with per-bar loss series persisted, then a
    battery MCS supersedes both the bounds and the tile MCS}
  \item The spectral-$k$NN headline remains a minimum over 32 cells selected on
    the evaluation panel. Bounds against the parametric arms
    (Section~\ref{sec:knn_selection}) do not address that selection.
  \item Everything in this section is conditional on the reported HAC standard
    errors being computed as documented. The bounds inherit whatever bandwidth
    convention produced the battery's $t$ statistics.
\end{itemize}
